\section{Introduction}

Reinforcement learning with verifiable rewards (\rlvr{}) has become one of the most effective and scalable paradigms for improving reasoning models. In mathematical reasoning, code generation, theorem proving, games, and other domains with deterministic checkers, a trajectory can be evaluated by an external rule rather than by a learned reward model. This property gives verifiable rewards a distinctive advantage: they are cheap to obtain at scale, robust to subjective preference drift, and nearly unbiased with respect to the task outcome. Unlike learned reward models, which may introduce evaluator bias or create reward-hacking incentives, a verifier answers a much simpler question: did the final object satisfy the task specification?

This advantage, however, comes with a severe cost. Verifiable rewards trade reward bias for feedback bandwidth. A complex trajectory may contain many decisions, intermediate beliefs, partial discoveries, and local errors, but the verifier typically compresses all of them into a single terminal scalar, often a binary $0$--$1$ signal. The reward is not weak because it is wrong; it is weak because it is narrow. From the perspective of policy optimization, the verifier is a low-bandwidth feedback channel: it preserves the correctness of the final outcome while discarding most information about how that outcome was produced.

This low-bandwidth view helps explain several puzzling phenomena in recent \rlvr{} systems. When all rollouts for a prompt receive the same reward, group-relative methods obtain little or no useful advantage signal, even though the prompt may still contain learnable structure. When a reward is only observed at the end of a long reasoning trace, credit assignment becomes ambiguous: the update does not know whether success or failure came from search strategy, concept selection, arithmetic execution, formatting, or accidental exploration. As the policy becomes more confident, rollout diversity can collapse, reducing the number of distinct outcomes exposed to the verifier. Conversely, methods that seem algorithmically different--dense rewards, process supervision, entropy shaping, adaptive rollout allocation, self-verification, and trajectory-level exploration--can often be understood as attempts to recover, redistribute, or better use the limited information carried by verifiable feedback \cite{hybrid_reinforcement_2026,exploration_exploitation_rlvr_2026,lookahead_rollouts_2026,eepo_2026,no_prompt_left_behind_2026,igpo_2026}.

The bottleneck is especially sharp in \llm{} \rlvr{} on a fixed prompt dataset. In classical reinforcement learning, interacting with an environment can continually expose new states and new reward events. In contrast, with a fixed dataset and deterministic verifier, the source of verifier information is strongly bounded by the prompt set, rollout budget, and reward channel. Training does not create new ground-truth feedback; it changes which parts of a finite feedback budget are revealed by sampling, how that feedback is allocated across prompts, and how efficiently it is converted into policy updates. Thus many \rlvr{} algorithms should be evaluated not only by final accuracy, but also by how much verifier information they expose, how strong an update they produce from that information, and how much useful learning signal is wasted through zero-variance groups, redundant trajectories, or overly sparse rewards.

This paper studies verifiable-reward learning through this information-bottleneck lens. Our goal is not to argue that verifiable rewards should be replaced by learned rewards. Instead, we separate two issues that are often conflated: reward bias and reward information capacity. Verifiable rewards are attractive because they reduce bias; their limitation is that they provide a narrow channel through which task information reaches the optimizer. Understanding this trade-off allows us to compare AlphaZero-like self-play, \llm{} reasoning, and diffusion generation under a common vocabulary.

We make the following contributions:
\begin{enumerate}
  \item We formulate verifiable-reward learning as a low-bandwidth feedback channel, making explicit the distinction between reward bias and reward information capacity.
  \item We introduce quantitative notions for dataset-level feedback budget, realized reward information, update strength, and information utilization. In particular, we emphasize that for \llm{} \rlvr{} on a fixed dataset with deterministic verifiers, the available verifier information is bounded and often effectively fixed by the dataset, rollout budget, and reward channel.
  \item We analyze the information bottleneck across three settings: AlphaZero-like self-play in board-game-style environments, \llm{} \rlvr{}, and diffusion-model \rlvr{}. Across these domains, we identify different ways to improve information efficiency, including richer rollout diversity, adaptive allocation of verifier queries, dense or trajectory-level feedback, and reward-aware credit assignment.
\end{enumerate}

By viewing verifiable rewards as an information bottleneck, we obtain a unified explanation for why sparse outcome feedback can be reliable yet inefficient, why some prompts contribute little to learning despite being meaningful, and why recent algorithmic improvements often work by increasing the effective bandwidth between the verifier and the policy optimizer.
